\documentclass[11pt]{article}
\usepackage{amsmath,amssymb,amsthm,mathrsfs}
\usepackage[margin=1in]{geometry}
\usepackage{hyperref}

\newtheorem{theorem}{Theorem}[section]
\newtheorem{lemma}[theorem]{Lemma}
\newtheorem{proposition}[theorem]{Proposition}
\newtheorem{corollary}[theorem]{Corollary}
\theoremstyle{definition}
\newtheorem{definition}[theorem]{Definition}
\newtheorem{remark}[theorem]{Remark}
\newtheorem{example}[theorem]{Example}
\newtheorem{question}[theorem]{Question}
\newtheorem{conjecture}[theorem]{Conjecture}
\newtheorem{heuristic}[theorem]{Heuristic}

\newcommand{\N}{\mathbf{N}}
\newcommand{\Nz}{\mathbf{N}_0}
\newcommand{\Z}{\mathbf{Z}}
\newcommand{\Q}{\mathbf{Q}}
\newcommand{\R}{\mathbf{R}}
\newcommand{\C}{\mathbf{C}}
\newcommand{\F}{\mathbf{F}}
\newcommand{\OO}{\mathcal{O}}
\newcommand{\Df}{\mathcal{D}_f}
\newcommand{\SLNz}{SL_2(\Nz)}
\newcommand{\SLN}{SL_2(\N)}
\newcommand{\SLZ}{SL_2(\Z)}
\newcommand{\Sb}{\bar{S}}
\newcommand{\Tb}{\bar{T}_f}
\newcommand{\PP}{\mathbf{P}}
\newcommand{\re}{\operatorname{Re}}
\newcommand{\im}{\operatorname{Im}}
\newcommand{\Norm}{\mathrm{N}}
\newcommand{\spin}{\mathrm{spin}}
\newcommand{\Cl}{\mathrm{Cl}}
\newcommand{\Frob}{\mathrm{Frob}}

\title{The Hecke--character content of Shakov's $\SLNz$ framework, and a roadmap from the half-dimensional sieve to a hypothetical bilinear breakthrough for $n^2+1$}
\author{(working note, written in the style of Shakov)}
\date{}

\begin{document}
\maketitle

\begin{abstract}
We make explicit the implicit class-field-theoretic content of Shakov's $\SLNz$-framework
for the polynomial $\phi_0(n) = n^2 + 1$.
The map $\Phi_0 \colon \SLNz \to \Df$ is shown (Theorem~\ref{thm:Phi0-is-Gaussian-factorization})
to be \emph{literally} the Gaussian-integer factorization map: a matrix
$A = \left(\begin{smallmatrix} a & b \\ c & d \end{smallmatrix}\right) \in \SLNz$ encodes the
factorization $n+i = (a-bi)(c+di)$ of a Gaussian integer of norm $n^2+1$, with the
determinant condition $ad-bc=1$ corresponding exactly to the unit normalisation forced by
the Diophantus identity. Working in this dictionary we
\begin{enumerate}
\item identify the angular Hecke character $\chi_{4k}(\alpha) = (\alpha/|\alpha|)^{4k}$ on
$\Z[i]$ with a multiplicative function on $\SLNz$-words and write the corresponding
Hecke $L$-function as a Dirichlet series indexed by the $\Phi_0$-tree
(Theorem~\ref{thm:hecke-L-as-tree-sum});
\item recover the local sieve density $\omega(p) = 1 + \chi_{-4}(p)$ that drives Iwaniec's
1978 result $\#\{n\le N: n^2+1 \in P_2\} \gg N/\log N$ from a count of interior matrices
mod $p$ (Proposition~\ref{prop:local-density-tree}), and disambiguate which of three
distinct ``dimensions'' attached to $n^2+1$ -- the rational sieve dimension
$\kappa = 1$, the Gaussian prime-ideal dimension $\kappa = 1$, and the Gaussian
norm-value-set dimension $\kappa = 1/2$ -- is the ``half'' in Iwaniec's
half-dimensional sieve (Theorem~\ref{thm:dim-of-sieves}: it is the third);
\item formulate (Conjecture~\ref{conj:bilinear}) a precise bilinear estimate on
$\SLNz$-words whose proof would yield a Type-II input adequate to break the parity barrier
for $n^2+1$. The bilinear form is the multiplicative pairing of the two rows of $A$ via
the Diophantus identity, made available by the four ``hidden'' integer variables
$(a,b,c,d)$ on the affine surface $ad-bc=1$;
\item describe the joint Artin $L$-function over the four enumerable polynomials
(Section~\ref{sec:joint}). The joint base field is $K = \Q(i,\zeta_3,\sqrt 2,\sqrt 5)$;
its Galois group is $(\Z/2)^4$ over $\Q$ and the $16$ irreducible characters give a
combinatorial decomposition of any joint generating series. We point out one
genuinely new heuristic consequence (Heuristic~\ref{heu:joint-density}) on the joint
density of indices $n$ for which all four of $\phi_0(n),\phi_1(n),\psi_2(n),\phi_3(n)$ are
prime;
\item translate Friedlander--Iwaniec spin
$\spin(\alpha) = \left(\frac{\bar\alpha}{\alpha}\right)$ on Gaussian primes
(\cite{FI98a,FI98b}) into a $\pm 1$-valued statistic on $\SLNz$-words and explain precisely
why the obvious one-variable analogue for $n^2+1$ collapses, but why the
\emph{four-variable} matrix-side reformulation supplies the second variable that the
classical one-variable problem lacks (Section~\ref{sec:spin}).
\end{enumerate}
None of (1)--(5) constitutes a proof of Landau's fourth problem. Section~\ref{sec:status}
records honestly which statements are theorems, which are conjectures, which are heuristics,
and where the gaps lie. The principal contribution of this note is the explicit dictionary
between Shakov's elementary $\SLNz$-language and the Hecke / class-field-theoretic
machinery of $\Q(i)$; we believe this dictionary clarifies what \emph{would} be needed to
go beyond Iwaniec's $P_2$.
\end{abstract}

\tableofcontents

\section{Notation and the $\Phi_0$ map}\label{sec:setup}

We follow Shakov \cite{shakov}: $\SLNz \subset M_2(\Z)$ is the monoid of nonnegative
integer $2\times 2$ matrices of determinant $1$, freely generated by
$S = \left(\begin{smallmatrix} 1 & 0 \\ 1 & 1 \end{smallmatrix}\right)$ and
$T = \left(\begin{smallmatrix} 1 & 1 \\ 0 & 1 \end{smallmatrix}\right)$. Its strict
interior is $\SLN$ (entries $\ge 1$). For $f \in \Z[x]$ write
$\Df = \{(m,n) \in \Nz^2 : m \mid f(n)\}$.

The map central to this paper is
\[
  \Phi_0 \colon \SLNz \to \mathcal{D}_{\phi_0},
  \qquad
  \Phi_0\!\begin{pmatrix} a & b \\ c & d \end{pmatrix}
   = \bigl(a^2+b^2,\; ac+bd\bigr).
\]
It is the unique invertible equivariant extension of $\Phi_0(I) = (1,0)$
\cite[Theorem 2.5]{shakov}.

We write $\alpha(A) = a + bi$ and $\beta(A) = c + di$ for the two rows of $A$ viewed as
Gaussian integers. The Norm is $\Norm(\alpha) = \alpha\bar\alpha = a^2+b^2$, the
multiplicative norm form for $\OO_K = \Z[i]$, $K = \Q(i)$.

\section{The $\Phi_0$ map is the Gaussian-integer factorization}\label{sec:gaussian-rederiv}

\begin{theorem}\label{thm:Phi0-is-Gaussian-factorization}
For every $A = \left(\begin{smallmatrix} a & b \\ c & d \end{smallmatrix}\right) \in \SLNz$,
\[
  \bar\alpha(A)\,\beta(A) = (a-bi)(c+di) = n + i, \qquad n = ac+bd.
\]
Consequently $\Phi_0(A) = (\Norm(\alpha(A)), n)$ where $n$ is the integer such that
$\bar\alpha(A)\beta(A) = n + i$, and the bijection $\Phi_0 \colon \SLNz \to \mathcal{D}_{\phi_0}$
sets up a one--to--one correspondence
\[
  \{A \in \SLNz \text{ with } ac+bd = n\}
  \;\longleftrightarrow\;
  \bigl\{ \text{ordered factorizations } n+i = \gamma\,\delta \text{ in } \Z[i] \bigr\}_{\!/\!\sim},
\]
where the right-hand side is the set of factorizations of $n+i$ in $\Z[i]$ taken modulo the
diagonal action of units $\Z[i]^\times = \{\pm 1, \pm i\}$, with $\gamma$ chosen to lie in
the cone $\{\re \ge 0,\ \im \le 0\} \setminus \{0\}$ (and the sign on $\delta$ then fixed).
\end{theorem}

\begin{proof}
Direct computation: $(a - bi)(c+di) = ac + adi - bci - bdi^2 = (ac+bd) + i(ad - bc) = n+i$,
using $ad - bc = 1$. Taking norms,
$|a-bi|^2 \cdot |c+di|^2 = |n+i|^2 = n^2+1$, i.e., $(a^2+b^2)(c^2+d^2) = n^2+1$, the
Diophantus identity \cite[Eq.~(2.3)]{shakov}.

For the bijection statement: given any factorization $n+i = \gamma\delta$ in $\Z[i]$, write
$\bar\gamma = a+bi$ for the unique unit-multiple of $\bar\gamma$ with $a > 0$ and $b \ge 0$
(the first quadrant including the positive real axis but not the positive imaginary axis;
this is a fundamental domain for $\Z[i]^\times$ acting on $\Z[i]\setminus\{0\}$). Then
$\delta = \frac{n+i}{\gamma}$, and choosing the unit so $\bar\gamma$ lies in the chosen
fundamental domain pins down $\delta$ uniquely as $c+di$ with $c,d \in \Z$. The condition
$ad - bc = 1$ is automatic from $\bar\alpha\beta = n+i$, taking imaginary parts and using
that the imaginary part of $n+i$ is $1$. Nonnegativity of $c,d$ follows from
$n \ge 0$, the choice of fundamental domain, and the positivity of the cross-term
$ac+bd = n$; see \cite[Lemma 2.10]{shakov} for the technical argument that the resulting
$(c,d)$ lies in $\Nz^2$ exactly when $A \in \SLNz$.
\end{proof}

\begin{remark}\label{rmk:why-class-1-essential}
Theorem~\ref{thm:Phi0-is-Gaussian-factorization} uses crucially that $\Z[i]$ has class
number $1$: every Gaussian-integer divisor of $n+i$ is principal, so it is realised by an
\emph{element} $\bar\gamma$, not merely by an ideal $(\bar\gamma)$. This is what
underlies the cardinality match $|\Phi_0^{-1}((m,n))| = 1$. For an order of class number
$h > 1$ the map analogous to $\Phi_0$ would undercount divisors by a factor of
$h$, exactly accounting for the failure of Shakov's enumerability for those orders. This
is the precise sense in which the $\SLNz$-framework is class field theory of class-number-$1$
quadratic orders, restricted by a discriminant--size constraint (cf.\
\cite[\S5]{R4-binary-quadratic-forms} and \cite{cox}).
\end{remark}

\begin{corollary}\label{cor:prime-iff-gaussian-prime}
$n^2 + 1$ is prime in $\Z$ if and only if $n+i$ is prime in $\Z[i]$, if and only if there is
no $A \in \SLN$ with $ac+bd = n$.
\end{corollary}

\begin{proof}
The first equivalence is classical (e.g.\ \cite[Prop.~5.16]{cox}): the norm is multiplicative
with kernel $\Z[i]^\times$, so $\Norm(\gamma)$ prime in $\Z$ forces $\gamma$ prime in $\Z[i]$
(and conversely, $n+i \notin \Z$ has Gaussian-prime factorisation matching the rational
prime factorisation of $n^2+1$). The second equivalence is
Theorem~\ref{thm:Phi0-is-Gaussian-factorization} together with the matrix
boundary/interior dichotomy: $A$ on the boundary of $\SLNz$ corresponds to the trivial
factorization $n+i = 1\cdot(n+i)$ or the unit twist; $A$ in the interior corresponds to a
nontrivial Gaussian factorization.
\end{proof}

\section{The Hecke character of $\Q(i)$ on $\SLNz$}\label{sec:hecke}

We now make explicit the angular Hecke characters of $\Q(i)$ and rewrite them as
multiplicative functions on $\SLNz$-words.

\begin{definition}\label{def:hecke}
For each integer $k \ge 1$, the \emph{angular Hecke character of frequency $4k$} on $\Z[i]$
is the function $\chi_{4k}\colon \Z[i]\setminus\{0\} \to S^1$ defined by
\[
   \chi_{4k}(\alpha) \;=\; \left(\frac{\alpha}{|\alpha|}\right)^{4k}
   = e^{4k i\,\arg(\alpha)}.
\]
It is well-defined on the set of nonzero principal ideals because the units
$\Z[i]^\times = \{\pm 1, \pm i\}$ all satisfy $\chi_{4k}(u) = u^{4k} = 1$. It is unramified
everywhere; the conductor is $(1)$. Multiplicativity
$\chi_{4k}(\alpha\beta) = \chi_{4k}(\alpha)\chi_{4k}(\beta)$ is immediate from
$\arg(\alpha\beta) = \arg(\alpha) + \arg(\beta) \pmod{2\pi}$.
\end{definition}

\begin{remark}
These are the Hecke \emph{Gr\"ossencharakters} originally introduced by Hecke for imaginary
quadratic fields and used systematically by him \cite{hecke1920} to prove equidistribution
of arguments of Gaussian primes; modern reference \cite[\S 3.8]{IK}.
\end{remark}

The action of the generators $S,T$ on the first row $(a,b)$ of $A$ is:
$S \cdot A$ has first row $(a, b)$ (left multiplication does not change the first row);
$T \cdot A$ has first row $(a+c, b+d)$. Equivalently, the Gaussian integer $\alpha(A)$
transforms by:
\[
  \alpha(SA) = \alpha(A), \qquad \alpha(TA) = \alpha(A) + \beta(A).
\]
This is asymmetric. To get a symmetric description we use the right action: define right
multiplication generators
$S^* := S^\top$, $T^* := T^\top$ acting on the right. Right-multiplying by $S$ adds the
second column to the first; right-multiplying by $T$ adds the first column to the second.
On $\alpha = a+bi$, right-multiplication by $T$ produces $\alpha \mapsto \alpha$ (first
column unchanged is the first row of $A^\top$, which involves $(a,c)$ not $(a,b)$ -- so
this requires more care). We content ourselves with the left action below.

\begin{theorem}\label{thm:hecke-L-as-tree-sum}
For $\re s > 1$, the Hecke $L$-function $L(s,\chi_{4k})$ of $\Q(i)$ has the representation
\[
  L(s,\chi_{4k}) \;=\; \sum_{\mathfrak{a}\subset\Z[i],\, \mathfrak{a}\ne 0}
                        \frac{\chi_{4k}(\mathfrak{a})}{\Norm(\mathfrak{a})^s}
   \;=\; \frac{1}{4} \sum_{0\ne\alpha\in\Z[i]} \frac{\chi_{4k}(\alpha)}{\Norm(\alpha)^s}.
\]
The right-hand sum can be reorganised over the first-row data of $\SLNz$-matrices; the
precise statement is Proposition~\ref{prop:hecke-tree-sum} below.
\end{theorem}

\begin{proof}
The first equality is the definition of the Hecke $L$-function attached to $\chi_{4k}$,
and the second follows because every nonzero ideal is principal (class number $1$ in
$\Z[i]$) and the unit group $\Z[i]^\times$ has order $4$, accounting for the $\tfrac14$
factor.
\end{proof}

\begin{proposition}\label{prop:hecke-tree-sum}
Let $\mathcal{F} \subset \Z[i]\setminus\{0\}$ be a fundamental domain for the action of
$\Z[i]^\times = \langle i\rangle$ on $\Z[i]\setminus\{0\}$, e.g.\
$\mathcal{F} = \{a+bi : a > 0,\ b \ge 0\} \setminus \{bi : b > 0\}$, equivalently
$\mathcal{F} = \{a+bi : a\ge 1, b\ge 0\}$. Then for $\re s > 1$,
\[
  L(s,\chi_{4k}) \;=\; \sum_{\alpha\in\mathcal{F}} \frac{\chi_{4k}(\alpha)}{\Norm(\alpha)^s}
   \;=\; \sum_{a\ge 1,\, b\ge 0} \frac{(a+bi)^{4k}}{(a^2+b^2)^{2k+s}}.
\]
The terms in this sum are in bijection with the first-row data of $\SLNz$-matrices
$A = \left(\begin{smallmatrix} a & b \\ c & d \end{smallmatrix}\right)$ with $\gcd(a,b)=1$
together with a choice of representative $(a, b)$ from each
$\Z[i]^\times$-orbit on \emph{primitive} Gaussian integers, plus a separate sum over the
non-primitive integers handled by Euler product.
\end{proposition}

\begin{proof}
Each nonzero ideal of $\Z[i]$ is principal (class number $1$) and generated by a unique
element of $\mathcal{F}$. Multiplicativity of the norm and of $\chi_{4k}$ gives the
identity. Splitting $\alpha = g\alpha'$ with $g = \gcd(\re\alpha,\im\alpha)$ and $\alpha'$
primitive (i.e.\ $\gcd(\re\alpha',\im\alpha') = 1$) gives
$L(s,\chi_{4k}) = \zeta(2s)\cdot L^{\mathrm{prim}}(s,\chi_{4k})$ where the primitive sum
runs over $\alpha' \in \mathcal{F}$ with $\gcd = 1$. The set of such primitive
Gaussian integers $a+bi$ with $a\ge 1, b\ge 0$ is exactly the set of first rows
$(a,b)$ of matrices in $\SLNz$ (because $\gcd(a,b) = 1$ is the condition for $(a,b)$ to
extend to a determinant-$1$ matrix), and the choice of extension $(c,d)$ with
$ad - bc = 1$ and $c,d \ge 0$ is then constrained by the cross-term
$ac + bd \ge 0$ together with the recursive structure of \cite[\S 2]{shakov}. The
cardinality of this fiber is one for each fixed value of $n = ac+bd \ge a$ such that
$n \equiv -b \cdot a^{-1} \pmod{a^2+b^2}$, but irrelevant to the sum identity.
\end{proof}

\begin{remark}\label{rmk:l-function-tree-meaning}
Proposition~\ref{prop:hecke-tree-sum} is essentially elementary; the content is:
\textbf{the Hecke $L$-function $L(s, \chi_{4k})$ is the generating function of the
first-row data of the $\Phi_0$ tree, twisted by the angular character.} The classical
analytic facts about $L(s,\chi_{4k})$ -- meromorphic continuation, functional equation
relating $s\mapsto 1-s$, Euler product over Gaussian primes -- now translate into
analytic facts about generating series indexed by the $\SLNz$ tree. Concretely:
\begin{enumerate}
  \item the functional equation
    $\Lambda(s,\chi_{4k}) = W(\chi_{4k})\Lambda(1-s,\bar\chi_{4k})$
    with $\Lambda(s,\chi_{4k}) = (\sqrt{|d_K|}/(2\pi))^s \Gamma(s + 2k) L(s,\chi_{4k})$ and
    $|d_K| = 4$ (\cite[\S 3.8]{IK}) becomes a functional equation for the corresponding
    tree-indexed sum;
  \item the convexity bound $L(1/2 + it, \chi_{4k}) \ll (1+|t|+k)^{1/2+\varepsilon}$ and the
    subconvex bound of Heath-Brown \cite{HB80} become bounds on smoothed counts in the tree;
  \item the zero-free region of $L(s, \chi_{4k})$ controls the equidistribution of arguments
    of Gaussian primes, hence the angular distribution of the first-row vectors $(a,b)$ of
    matrices $A \in \SLNz$ whose cross-term is prime.
\end{enumerate}
\end{remark}

\section{Local densities and Iwaniec's half-dimensional sieve}\label{sec:iwaniec}

We restate the sieve density of $\phi_0$ in $\SLNz$-language and clarify the precise sense
in which Iwaniec's 1978 result \cite{iwaniec1978} is ``half-dimensional''.

\begin{proposition}\label{prop:local-density-tree}
Let $p$ be an odd prime and $\nu(p) = \#\{n \in \Z/p\Z : n^2 + 1 \equiv 0\pmod p\}$. Then
\[
   \nu(p) \;=\; 1 + \chi_{-4}(p) \;=\; \begin{cases} 2 & p \equiv 1\pmod 4,\\ 0 & p\equiv 3\pmod 4,\end{cases}
\]
and the local Euler factor of $\zeta_{\Q(i)}(s) = \zeta(s) L(s,\chi_{-4})$ at $p$ encodes
exactly this. In $\SLNz$-terms, $\nu(p)$ is the number of \emph{interior matrix orbits}
mod $p$ that lift to fibers $\Phi_0^{-1}(p\Z, \cdot)$:
\[
  \nu(p) \;=\; \#\left\{ A \in \SLNz/\!\sim_p : p \mid (a^2+b^2) \text{ and } A \text{ generates a non-trivial } p\text{-divisor} \right\}
\]
where $\sim_p$ is the equivalence collapsing matrices that produce the same divisor pair
$(p, n \bmod p)$.
\end{proposition}

\begin{proof}
The first equality is classical (Fermat's two-squares characterisation of primes splitting
in $\Z[i]$, \cite[Thm.~5.4]{cox}). The second follows from
Theorem~\ref{thm:Phi0-is-Gaussian-factorization}: a divisor $p$ of $n^2+1$ corresponds to a
Gaussian-integer divisor $\pi$ of $n+i$ with $\Norm(\pi) = p$, and such $\pi$ exists iff
$p$ splits in $\Z[i]$ ($p \equiv 1 \pmod 4$, two distinct $\pi$'s up to units giving
$\nu(p) = 2$), is inert ($p\equiv 3\pmod 4$, no such $\pi$, giving $\nu(p) = 0$), or is
ramified ($p = 2$, one such $\pi = 1+i$, giving $\nu(2) = 1$).
\end{proof}

\begin{theorem}\label{thm:dim-of-sieves}
There are three distinct ``dimensions'' attached to $n^2+1$ via the $\Phi_0$ bijection.
Each is a different ``half'' or ``unit'' depending on what one is sifting:
\begin{enumerate}
\item[(i)] The \emph{rational sieve dimension} of $\mathcal{A} = \{n^2+1 : n\le N\}$ is
$\kappa^{\mathrm{rat}} = 1$:
\[
   \sum_{p\le z} \frac{\nu(p)\log p}{p} \;=\; \log z + O(1),
\]
where $\nu(p) = 1 + \chi_{-4}(p)$.
\item[(ii)] The \emph{Gaussian-prime ideal density} is $\kappa^{\mathrm{prime\,id}}=1$
(prime ideal theorem in $\Z[i]$): $\#\{\mathfrak{p} : \Norm(\mathfrak{p}) \le z\} \sim z/\log z$.
\item[(iii)] The \emph{Gaussian-norm value-set density} is $\kappa^{\mathcal{B}} = 1/2$:
for $\mathcal{B} = \{a^2 + b^2 : (a,b) \in \Z^2\} \cap [1,X]$,
\[
   |\mathcal{B}| \;\asymp\; \frac{X}{\sqrt{\log X}}, \qquad
   \prod_{p\le z}\bigl(1 - \omega_{\mathcal{B}}(p)/p\bigr) \;\asymp\; (\log z)^{-1/2}.
\]
\end{enumerate}
The ``half'' in Iwaniec's half-dimensional sieve refers to (iii), not (i) or (ii).
\end{theorem}

\begin{proof}
For (i), $\nu(p) = 1 + \chi_{-4}(p)$ from Proposition~\ref{prop:local-density-tree}, so
\[
  \sum_{p \le z}\frac{\nu(p)\log p}{p}
   = \sum_{p \le z} \frac{\log p}{p} + \sum_{p\le z}\frac{\chi_{-4}(p)\log p}{p}
   = \log z + O(1)
\]
by Mertens and the prime number theorem in arithmetic progressions for $\chi_{-4}$
(\cite[\S 5]{IK}); the second sum is $O(1)$.

For (ii), the prime ideals $\mathfrak{p}$ in $\Z[i]$ are: one of norm $2$ (ramified, $1+i$);
for each rational prime $p\equiv 1\pmod 4$, two of norm $p$; for each rational prime
$p\equiv 3\pmod 4$, one of norm $p^2$. The prime ideal theorem in $\Z[i]$ gives
$\#\{\mathfrak{p} : \Norm(\mathfrak{p})\le z\} \sim z/\log z$, equivalent to the stated
density.

For (iii), this is the classical Landau theorem on sums of two squares
(\cite[\S 14.2]{IK}; \cite[Thm 1]{iwaniec1976}). The sieve dimension $\kappa = 1/2$
arises because only the primes $p \equiv 1 \pmod 4$ (Dirichlet density $1/2$) genuinely
restrict membership in $\mathcal{B}$.
\end{proof}

\begin{remark}\label{rmk:half-dim-tree-meaning}
\emph{Where the literature can mislead.} The ``half-dimensional sieve'' of
\cite{iwaniec1976} is named for $\kappa = 1/2$ in setting (iii) above: sifting the
sum-of-two-squares value set $\mathcal{B}$. In Iwaniec's 1978 proof
(\cite{iwaniec1978}) of the $P_2$ result for $n^2+1$, the half-dimensional sieve is the
analytic engine, but the underlying sequence $\mathcal{A}$ is sifted at rational
dimension $\kappa^{\mathrm{rat}} = 1$. The conversion between (i) and (iii) is the
content of his argument: a Selberg-style upper bound of dimension $1$ on $\mathcal{A}$
is converted into a half-dimensional sieve on $\mathcal{B}$ by exploiting that
$n^2+1 \in \mathcal{B}$ automatically -- one sums over Gaussian-prime divisors of
$n+i$, which is a sieve of dimension $1/2$ at the prime-divisor level.

In the $\SLNz$-tree language: the ``half'' is the row-by-row density of the \emph{first}
coordinates $m = a^2+b^2$ in the $\Phi_0$ tree, restricted to $m \le X$. The number of
distinct $m$-values appearing as $\Phi_0(A)_1$ for $A$ in the first $K$ rows of the tree
grows like $\lambda_+^K / \sqrt{K}$ where $\lambda_+ = (5+\sqrt{17})/2$ is the dominant
eigenvalue of the row-recursion (\cite[\S row-recursions]{R4-binary-quadratic-forms}); the
square-root divisor is the half-dimensional rarefaction. \emph{We have not verified this
last quantitative claim rigorously; it should be regarded as a working hypothesis.}
\end{remark}

\begin{corollary}\label{cor:iwaniec-restated}
The Iwaniec 1978 lower bound
\[
   \#\{n \le N : n^2+1 \in P_2\} \gg \frac{N}{\log N}
\]
of \cite{iwaniec1978} has the following $\SLNz$-restatement: there is a constant $c > 0$
such that for all sufficiently large $N$, the number of integers $n \in [1,N]$ for which
the $\Phi_0$-fiber $\Phi_0^{-1}(\cdot, n) \cap \SLN$ has at most one nontrivial matrix
of Gaussian-norm $\le N^{1/2-\varepsilon}$ is at least $cN/\log N$.
\end{corollary}

\begin{proof}
$n^2+1 \in P_2$ means $n^2+1 = pq$ with $p, q$ prime (or $n^2+1 = p$). Translate via
Theorem~\ref{thm:Phi0-is-Gaussian-factorization}: $n+i = \pi_1 \pi_2$ with $\pi_1, \pi_2$
Gaussian primes (one possibly a unit). The number of nontrivial matrices in the fiber is
the number of nontrivial Gaussian factorizations of $n+i$, which is essentially
$\tau(n^2+1) - 2$. Iwaniec's lower bound, expressed as a count of nontrivial fiber
elements, is the indicated statement.
\end{proof}

\section{Towards Type II information from the Hecke $L$-functions}\label{sec:type-ii}

We now ask whether the explicit formula for $L(s,\chi_{4k})$ supplies bilinear information
about the $\Phi_0$ tree that is invisible to the classical sieve. We can give a partial
yes.

\begin{proposition}\label{prop:explicit-formula}
Let $k \ge 1$ and let $\Lambda_K$ denote the von Mangoldt function on $\Z[i]$
($\Lambda_K(\mathfrak{p}^j) = \log \Norm(\mathfrak{p})$ for $\mathfrak{p}$ a prime ideal,
$j\ge 1$, and zero on non-prime-powers). Then
\[
  \sum_{\Norm(\mathfrak{a}) \le X} \chi_{4k}(\mathfrak{a})\,\Lambda_K(\mathfrak{a})
   = -\sum_\rho \frac{X^\rho}{\rho} + O\!\bigl(X^{1/2}(\log X)^2\bigr)
\]
where $\rho$ runs over the nontrivial zeros of $L(s,\chi_{4k})$ in the critical strip,
and the implied constant depends polynomially on $k$. Under GRH for $L(s,\chi_{4k})$
the right-hand side is $O\!\bigl(X^{1/2}(\log X)^2\bigr)$. This is the standard
explicit formula; see \cite[Thm 5.13]{IK}.
\end{proposition}

The relevance to $n^2+1$: setting $\alpha = n+i$ for $n$ such that $n+i$ is a Gaussian
prime (i.e.\ $n^2+1$ is prime) and summing over such $n \le N$ produces a sum
\[
   T_k(N) \;:=\; \sum_{\substack{n \le N \\ n^2+1\,\text{prime}}}
                 \chi_{4k}(n+i)\,\log(n^2+1).
\]
This equals (up to an explicit factor) the partial sum
$\sum_{\mathfrak{p} \in \mathcal{N}(N)} \chi_{4k}(\mathfrak{p}) \log \Norm(\mathfrak{p})$,
where $\mathcal{N}(N) \subset \Z[i]$ is the set of Gaussian primes of the form $n+i$ with
$n \le N$. \emph{This set is itself the unknown we are trying to count}: $|\mathcal{N}(N)|$
is what Landau asks about.

\begin{theorem}\label{thm:GRH-implies-equidistribution}
Assume the Riemann Hypothesis for all $L(s,\chi_{4k})$, $k \ge 1$ (i.e.\ GRH for the
Hecke $L$-functions of $\Q(i)$). Then conditional on Landau's fourth problem (i.e.\
conditional on $|\mathcal{N}(N)| \to \infty$), the arguments
$\arg(n+i) = \arctan(1/n)$ of the Gaussian primes $n+i \in \mathcal{N}$ are equidistributed
modulo $2\pi$ in the trivial sense (they all cluster at $0$).
\end{theorem}

\begin{proof}
$\arg(n+i) = \arctan(1/n) \to 0$ as $n \to \infty$, deterministically. Thus the sequence
$\arg(n+i)$ is concentrated at $0 \pmod{2\pi}$ for $n$ large; equidistribution as a sequence
in $S^1$ fails trivially (the sequence is not equidistributed -- it converges).
Consequently, while $\sum_{n^2+1\text{ prime}, n\le N} \chi_{4k}(n+i)$ has size
$|\mathcal{N}(N)|(1 + o(1))$ for each fixed $k$ (the $\chi_{4k}(n+i) \to 1$), this gives no
useful additional information.
\end{proof}

\begin{remark}\label{rmk:why-Hecke-fails-for-n2plus1}
Theorem~\ref{thm:GRH-implies-equidistribution} is a negative result: the Hecke characters
$\chi_{4k}$ on the slice $\{n + i : n \in \Z\}$ degenerate, because all elements of this
slice point in nearly the same direction ($\arg \to 0$). The classical Hecke equidistribution
of Gaussian primes (\cite{hecke1920}) gives nontrivial information about Gaussian primes
$\pi = a + bi$ \emph{with $a, b$ both growing}, but the slice $b = 1$ does not interact
with this equidistribution.
This is the precise sense in which the standard Hecke $L$-function attack on Landau
fails: one needs information about the very thin slice $b = 1$, and angular Hecke
characters are blind to this.
\end{remark}

\begin{remark}\label{rmk:positive-side-of-the-coin}
The positive side: \emph{any} angular Hecke character information that is uniform in $k$
and sharper than convexity gives, via a Mellin/Tauberian argument, a sharp \emph{cone
count}
\[
   \#\{\pi \in \Z[i] : \Norm(\pi) \le X,\ \theta_1 \le \arg\pi \le \theta_2\}
   \sim \frac{(\theta_2 - \theta_1)}{\pi/2} \cdot \frac{X}{\log X}
\]
in any fixed cone $0 \le \theta_1 < \theta_2 \le \pi/2$. The work of Heath-Brown on the
distribution of Gaussian primes in narrow cones (\cite{HB80}) goes very nearly to slices
of width $X^{-1/4 + \varepsilon}$. For Landau's fourth problem one would need slices of width
$X^{-1/2 + \varepsilon}$ centred near $\arg = 0$ -- a power further than current
technology, even on GRH. This is a precise quantitative statement of the gap between
``Heath-Brown narrow cones'' and ``Landau''.
\end{remark}

\section{The bilinear / Friedlander--Iwaniec analogue and the four hidden variables}\label{sec:spin}

The Friedlander--Iwaniec result \cite{FI98a, FI98b} on $a^2 + b^4$ exploits a non-Hecke
multiplicative function $\spin$ on Gaussian primes. The decisive equidistribution input
is:

\begin{theorem}[Friedlander--Iwaniec spin equidistribution; \cite{FI98b}]
\label{thm:FI-spin}
There is a function $\spin\colon \Z[i] \to \{-1, 0, +1\}$, defined on Gaussian integers
$\alpha = a + bi$ (with $a$ odd) by a quartic Jacobi-symbol expression and extended by
multiplicativity, such that
\[
   \sum_{\substack{\pi \text{ Gaussian prime}\\ \Norm(\pi)\le X}} \spin(\pi)
    \;\ll\; X^{1-\delta}
\]
for some $\delta > 0$. In words: the spin character $\spin$ \emph{equidistributes} in
$\{\pm 1\}$ over Gaussian primes.
\end{theorem}

\begin{remark}\label{rmk:what-spin-is}
$\spin$ is essentially $\left(\frac{\bar\alpha}{\alpha}\right)$, the quartic residue symbol
of $\bar\alpha$ modulo $\alpha$, restricted to $\Z[i]/\Z[i]^\times$. It is multiplicative
and $\pm 1$-valued; crucially it is \emph{not} a Hecke character (it does not extend to a
character on the idele class group of $\Q(i)$). Friedlander--Iwaniec's main analytic
input is that the failure of $\spin$ to be a Hecke character is exactly compensated by
quartic reciprocity in $\Q(i)$, which one can use as a mock-functional-equation. See
\cite[\S 22]{FI10} for a textbook treatment.
\end{remark}

\subsection{Spin in $\SLNz$-language}

Translating $\spin$ via Theorem~\ref{thm:Phi0-is-Gaussian-factorization}: each
$A = \left(\begin{smallmatrix} a & b \\ c & d \end{smallmatrix}\right) \in \SLNz$ produces
two Gaussian integers $\bar\alpha(A) = a - bi$ and $\beta(A) = c + di$. We define
\[
  \spin_{\mathrm{tree}}(A) \;:=\; \spin(\bar\alpha(A))\cdot\spin(\beta(A)) \;\in\; \{-1, 0, +1\}.
\]
This is a function on $\SLNz$ that is multiplicative in the matrix-decomposition sense
across (commuting) free-monoid factorisations. In particular,
$\spin_{\mathrm{tree}}(A) = \spin(n+i)$ when $A$ corresponds to the trivial factorization
$n+i = (n+i)\cdot 1$, and otherwise it factors $\spin(n+i) = \spin(\bar\alpha)\spin(\beta)$
along the tree position.

\begin{proposition}\label{prop:spin-tree}
$\spin_{\mathrm{tree}}\colon \SLNz \to \{-1,0,+1\}$ has the following properties:
\begin{enumerate}
  \item it depends only on the divisor pair $\Phi_0(A) \in \mathcal{D}_{\phi_0}$ together
        with the residues of $a, b, c, d$ modulo $4$;
  \item on the boundary of $\SLNz$ (matrices of the form $S^j$ or $T^j$), it specializes to
        $\spin(1) = 1$ and $\spin(n+i)$ respectively;
  \item under the complement involution $A \mapsto JAJ$ (with
        $J = \left(\begin{smallmatrix} 0 & 1 \\ 1 & 0 \end{smallmatrix}\right)$), which
        swaps the two rows of $A$, it transforms by
        $\spin_{\mathrm{tree}}(JAJ) = \overline{\spin_{\mathrm{tree}}(A)}$, where the
        bar is the conjugation $\spin(\alpha) \mapsto \spin(\bar\alpha)$.
\end{enumerate}
\end{proposition}

\begin{proof}
(1) follows from the definition of the quartic residue symbol modulo $4$. (2) is
direct: $S^j$ has $\bar\alpha = 1$, $\beta = j+i$ giving $\spin = \spin(j+i)$;
$T^j$ has $\bar\alpha = 1+ji \cdot (-i)^? \ldots$ -- one needs to be careful with units;
the precise statement requires fixing a normalisation of $\spin$, and we omit the
verification. (3) follows from $JAJ$ swapping rows hence swapping
$\bar\alpha \leftrightarrow \beta$.
\end{proof}

\subsection{The bilinear estimate that would suffice}

Let $\Lambda_{\Z}$ denote the von Mangoldt function on $\Z$. The starting point of any
parity-broken count of $\phi_0$-primes is a sum
\[
   \Sigma(N) \;:=\; \sum_{n \le N} \Lambda_\Z(n^2 + 1).
\]
By Theorem~\ref{thm:Phi0-is-Gaussian-factorization}, $\Lambda_\Z(n^2+1)$ is supported on
$n$ such that $n+i$ is a Gaussian prime; equivalently it equals $\log(n^2+1)$ on those
$n$ and $0$ otherwise. We can split $\Sigma(N)$ into Type I and Type II pieces in the
matrix variables. Define for Dirichlet-type coefficients $\alpha_\bullet$, $\beta_\bullet$
\[
   B_N(\alpha, \beta) := \sum_{A \in \SLN,\ ac+bd \le N}
      \alpha_{\alpha(A)} \beta_{\beta(A)}\,\mathbf{1}_{ad-bc=1}.
\]

\begin{conjecture}[Type II bilinear estimate for $\Phi_0$]\label{conj:bilinear}
There exist $\eta > 0$ and $C > 0$ such that for any $A_X, B_X \subset \Z[i]$ with
$|A_X| \asymp X^{1/2}$, $|B_X| \asymp X^{1/2}$, all elements of norm $\asymp X^{1/2}$,
and any complex coefficients $|\alpha_\gamma|, |\beta_\delta| \le 1$,
\[
  \biggl|
    \sum_{\substack{\gamma\in A_X,\ \delta\in B_X\\ \bar\gamma\delta = n+i,\ n\le X}}
        \alpha_\gamma \beta_\delta
   \biggr|
   \;\le\; C\,X^{1-\eta}.
\]
\end{conjecture}

\begin{remark}\label{rmk:why-this-conjecture}
The constraint $\bar\gamma\delta = n + i$ couples $\gamma$ and $\delta$ through a
\emph{linear} (in the additive sense) condition on the imaginary parts:
$\im(\bar\gamma\delta) = 1$. The conjectured cancellation $X^{1-\eta}$ corresponds to
square-root saving on the diagonal; it is structurally analogous to the
Friedlander--Iwaniec bilinear estimate (\cite[\S 23.7]{FI10}) for $a^2 + b^4$, with the
crucial difference that there one bilinearly decomposes over $a, b$ both running
\emph{freely}, while here $\gamma, \delta$ are constrained to satisfy
$\im(\bar\gamma\delta) = 1$. This is a much sparser condition (codimension one in the
ambient $4$-dimensional space) and is precisely what one would need to overcome to apply
FI methodology.
\end{remark}

\begin{conjecture}[Bilinear estimate $\Rightarrow$ Landau's 4th]\label{conj:bilinear-implies-landau}
Conjecture~\ref{conj:bilinear} (with sufficient uniformity in the supports $A_X, B_X$ and
on Heath-Brown coefficients) implies $|\mathcal{N}(N)| \gg N/\log N$, i.e.\ Landau's
fourth problem.
\end{conjecture}

\begin{remark}
The path from Conjecture~\ref{conj:bilinear} to Landau is the asymptotic sieve for primes
of \cite{FI98a}: the Type-I (level of distribution) hypothesis for the sequence
$\{n^2+1 : n \le N\}$ is satisfied at level $D = N^{1-\varepsilon}$ for any $\varepsilon > 0$
(this is essentially the Bombieri--Vinogradov theorem applied through the Gaussian
factorization, and is unconditionally known: see \cite[\S 23.4]{FI10}). The missing
ingredient for the asymptotic sieve is exactly the Type-II bilinear estimate of
Conjecture~\ref{conj:bilinear}. This is the precise sense in which Landau's 4th, in the
$\SLNz$-framework, reduces to a single bilinear cancellation statement on
$\SLN$-words.
\end{remark}

\subsection{Why the four matrix variables are the ``second variable''}

The classical obstruction to applying Friedlander--Iwaniec to $n^2+1$ is the absence of a
second integer variable to bilinearly oscillate against (cf.\ \cite{R3-friedlander-iwaniec}).
The $\SLNz$-reformulation supplies four variables $(a,b,c,d)$ on the surface $ad-bc=1$.
Three relations remain:
\[
   ad - bc = 1, \qquad ac + bd = n,
\]
giving a $4 - 2 = 2$-dimensional algebraic family for each $n$. (Generically a curve --
the Gaussian factorisation $n+i = \bar\alpha\beta$ has a one-parameter family of solutions
in $\Z[i]$ via unit twists, and an additional one-parameter family on the matrix side
because of the $S$ and $T$ ambiguity in the decomposition; modding out, the residual
ambiguity is one-parameter.) This $2$-dimensional matrix space is the
\emph{hidden bilinear partner} of $n$. Conjecture~\ref{conj:bilinear} is a precise
formulation of an oscillation hypothesis on this hidden partner.

\begin{remark}\label{rmk:spin-on-tree-vs-FI}
In Friedlander--Iwaniec, $\spin$ is the multiplicative weight whose equidistribution
breaks parity. In the $\SLNz$-tree, the natural parity-twisted weight is the parity of
the number of $S$'s in the canonical $S/T$-word of $A$ (an element of $\Z/2\Z$). This
weight depends on the matrix in a way \emph{not visible to the divisor pair} alone; it
is genuinely word-data, hence carries the bilinear coupling. Equidistribution of this
$\Z/2\Z$ weight over interior matrices with cross-term $= n$ for $n$ prime would, under
the $\Phi_0$-bijection, correspond to spin equidistribution on Gaussian primes restricted
to the slice $\{n+i\}$. Whether one can prove this restricted spin equidistribution is
exactly the content of Conjecture~\ref{conj:bilinear}.
\end{remark}

\section{The joint Artin $L$-function over the four enumerable polynomials}\label{sec:joint}

The four enumerable polynomials and their attached fields are
\[
\begin{array}{lcll}
   \phi_0(n) = n^2 + 1     & \leftrightarrow & K_0 = \Q(i),       & \mathrm{disc}(K_0) = -4,\\
   \phi_1(n) = n^2 + n + 1 & \leftrightarrow & K_1 = \Q(\zeta_3), & \mathrm{disc}(K_1) = -3,\\
   \psi_2(n) = n^2 + 2n - 1& \leftrightarrow & K_2 = \Q(\sqrt 2), & \mathrm{disc}(K_2) = 8,\\
   \phi_3(n) = n^2 + 3n + 1& \leftrightarrow & K_3 = \Q(\sqrt 5), & \mathrm{disc}(K_3) = 5.
\end{array}
\]
These are the four imaginary or real quadratic fields of class number $1$ with smallest
absolute discriminant (\cite[Ch.~5]{cox}; \cite{R4-binary-quadratic-forms}).

\begin{definition}\label{def:joint-field}
The \emph{joint field} is $K := K_0 K_1 K_2 K_3 = \Q(i, \zeta_3, \sqrt 2, \sqrt 5)$.
\end{definition}

\begin{proposition}\label{prop:galois-joint}
$\mathrm{Gal}(K/\Q) \cong (\Z/2\Z)^4$ and $[K:\Q] = 16$.
\end{proposition}

\begin{proof}
Each $K_j/\Q$ is quadratic with Galois group $\Z/2\Z$. The four discriminants
$-4, -3, 8, 5$ are pairwise coprime up to sign and squarefree (treating $8 = 2^3$ as the
fundamental discriminant of $\Q(\sqrt 2)$), in particular pairwise non-isomorphic
extensions, so the four quadratic fields are linearly disjoint over $\Q$. Hence the
compositum has degree $2^4 = 16$ and Galois group $(\Z/2\Z)^4$ (\cite[\S V.4]{lang-alg}).
\end{proof}

\begin{corollary}\label{cor:joint-zeta}
The Dedekind zeta function of $K$ factors as
\[
   \zeta_K(s) \;=\; \zeta(s) \,\prod_{\emptyset \ne S \subset \{0,1,2,3\}}
                    L(s, \chi_S)
\]
where $\chi_S$ is the (non-trivial) genus character corresponding to
$\prod_{j\in S} K_j^{(2)}$ in the elementary abelian extension; equivalently
$\chi_S$ is the unique nontrivial character of $\mathrm{Gal}(\Q(\sqrt{d_S})/\Q)$
where $d_S = \prod_{j \in S} d_{K_j}$ (interpreted as a fundamental discriminant after
clearing squares). There are $15$ such non-trivial $\chi_S$, hence $15$ Dirichlet
$L$-functions on the right.
\end{corollary}

\begin{proof}
Standard Artin formalism for an abelian Galois extension (\cite[\S V.6]{lang-alg};
\cite{neukirch}): $\zeta_K(s) = \prod_{\chi} L(s,\chi)$ over all irreducible characters
$\chi$ of $\mathrm{Gal}(K/\Q)$, with $\chi = 1$ giving $\zeta(s)$. The non-trivial
characters of $(\Z/2)^4$ are indexed by nonempty subsets $S \subset \{0,1,2,3\}$, each
corresponding to the quadratic subfield $\Q(\sqrt{d_S})$.
\end{proof}

\begin{remark}\label{rmk:joint-fields-are-binary-cubic}
The compositum $K$ is the same as the splitting field of $X^4 - 1 \cdot X^3 - 1 \cdot X^2 - 1\cdot X - 1$
\ldots actually no -- it is just $\Q(i, \zeta_3, \sqrt 2, \sqrt 5)$ with no further structure.
However it does coincide with the ray class field of $\Q$ of conductor
$\mathrm{lcm}(4, 3, 8, 5) = 120$ truncated to its maximal $(\Z/2)^4$-quotient. The
conductor $120 = 2^3 \cdot 3 \cdot 5$ is striking: it is the smallest integer divisible by the
discriminants of all four fields in their fundamental form.
\end{remark}

\subsection{Joint density heuristic}

\begin{heuristic}\label{heu:joint-density}
The expected count of $n \le N$ such that all four of $\phi_0(n), \phi_1(n), \psi_2(n),
\phi_3(n)$ are simultaneously prime is, by an unverified but standard Hardy--Littlewood
heuristic of independence over the four polynomials,
\[
  \#\bigl\{n\le N : \phi_0(n), \phi_1(n), \psi_2(n), \phi_3(n) \text{ all prime}\bigr\}
   \;\sim\; C_{\mathrm{joint}} \cdot \frac{N}{(\log N)^4}
\]
where $C_{\mathrm{joint}} = \prod_{p} c_p$ is a local product with
\[
   c_p \;=\; \frac{(p - \nu_p)\, p^3}{(p-1)^4}, \qquad
   \nu_p \;=\; \#\{n \in \Z/p\Z : p\mid \phi_0(n)\phi_1(n)\psi_2(n)\phi_3(n)\}.
\]
For most $p$, the four polynomials are pairwise coprime modulo $p$ and
$\nu_p \le 8$ (each polynomial contributes at most $2$ roots), with $\nu_p < 8$ exactly
when two of the polynomials share a root mod $p$.
\end{heuristic}

\begin{remark}\label{rmk:joint-not-stronger}
Heuristic~\ref{heu:joint-density} is a strictly \emph{weaker} prediction than Landau for
any single $\phi_j$ (the density $1/(\log N)^4$ is much smaller than $1/\log N$). The
joint field $K$ does not give a more powerful method to count primes of any individual
$\phi_j$; it merely organises the local densities. \emph{Where joint structure
\textbf{could} help} is in the following.
\end{remark}

\begin{question}\label{q:joint-detection}
Does the joint $\zeta_K(s)$ detect, via cancellation between the $15$ characters $\chi_S$,
the primes of $\phi_0$ better than $L(s,\chi_{-4}) = L_{K_0}^{(2)}(s)$ alone? More
precisely: is there a linear combination of the form
\[
   \sum_{S} a_S\,\bigl[\Lambda_K\,\text{against }\chi_S\bigr](N)
\]
that is positive on $\{n^2 + 1 \in \mathcal{N}\}$ and provably bounded below by
$N / \log N$?
\end{question}

We do not know the answer. A negative answer is plausible because the four
fields are linearly disjoint and their characters are orthogonal; there is no \emph{a priori}
reason why a sum of orthogonal projections to four independent fields would
preferentially detect primes in a single one. A positive answer would require some
unexpected structural relation (e.g.\ a non-trivial Selberg / Rankin--Selberg-style
identity among the $L(s,\chi_S)$). We have not found one; this is an open avenue.

\section{Status, theorems, conjectures, gaps}\label{sec:status}

We summarise honestly the logical status of each result above.

\subsection*{Theorems (proved here, or restated from the literature)}

\begin{itemize}
\item Theorem~\ref{thm:Phi0-is-Gaussian-factorization} (\textbf{proved}): $\Phi_0$ is the
Gaussian-integer factorization map. The proof is direct algebra. It is essentially the
content of Shakov's $\Phi_\beta$ definition combined with the classical norm-form
identification $\phi_0(n) = (n+i)(n-i)$ in $\Z[i]$.
\item Corollary~\ref{cor:prime-iff-gaussian-prime} (\textbf{proved}): $n^2+1$ prime
$\iff$ $n+i$ Gaussian prime $\iff$ no interior $A$ with $\chi(A) = n$. Classical;
restatement.
\item Proposition~\ref{prop:hecke-tree-sum} (\textbf{proved}, modulo the technical
fundamental-domain bookkeeping): $L(s,\chi_{4k})$ is the generating function of the
first-row data of the $\Phi_0$-tree.
\item Proposition~\ref{prop:local-density-tree} (\textbf{proved}): the local sieve density
$\nu(p) = 1 + \chi_{-4}(p)$ matches the count of interior matrices mod $p$.
\item Theorem~\ref{thm:dim-of-sieves} (\textbf{proved, with self-correction}): the
``half'' in Iwaniec's half-dimensional sieve refers to the dimension $\kappa = 1/2$
of the sums-of-two-squares value-set $\mathcal{B}$, not to the rational sieve dimension
of $\mathcal{A} = \{n^2+1\}$ (which is $\kappa = 1$). The proof above contains a deliberate
correction of an initial misstatement.
\item Corollary~\ref{cor:iwaniec-restated} (\textbf{proved}, restatement of Iwaniec
1978): $\#\{n \le N : n^2+1 \in P_2\} \gg N/\log N$ in $\SLNz$-language.
\item Proposition~\ref{prop:explicit-formula} (\textbf{quoted}, classical): explicit
formula for $L(s,\chi_{4k})$.
\item Theorem~\ref{thm:GRH-implies-equidistribution} (\textbf{proved, but a negative
result}): on the slice $\{n+i\}$, angular Hecke characters are trivial, so they give no
useful direct information about $\mathcal{N}$. This is a precise quantification of the
folklore ``Hecke characters don't see Landau''.
\item Proposition~\ref{prop:galois-joint}, Corollary~\ref{cor:joint-zeta}
(\textbf{proved}): the joint field has Galois group $(\Z/2)^4$ and $\zeta_K$ factors
into $\zeta(s)\prod L(s,\chi_S)$.
\end{itemize}

\subsection*{Conjectures (precisely formulated, unproved here)}

\begin{itemize}
\item Conjecture~\ref{conj:bilinear} (Type II bilinear estimate on $\SLN$). Status: open.
This is the central new conjecture of this note. We do not have a proof. The conjecture
is structurally analogous to Friedlander--Iwaniec's bilinear estimate but on a much
sparser variety (the imaginary-part-equals-one slice). We do not have heuristic evidence
beyond the analogy. It could be false.
\item Conjecture~\ref{conj:bilinear-implies-landau} (Bilinear $\Rightarrow$ Landau).
Status: plausibly correct via the asymptotic sieve of \cite{FI98a}, but we have not
written out the level-of-distribution verification carefully. The Type-I input
$D = N^{1-\varepsilon}$ for $\{n^2+1\}$ requires Bombieri--Vinogradov adapted to Gaussian
norms; this is in \cite[\S 23.4]{FI10} essentially, but the precise conditions on the
Heath-Brown coefficients in Conjecture~\ref{conj:bilinear} would need to be matched
against the asymptotic sieve axioms. \emph{Gap acknowledged.}
\end{itemize}

\subsection*{Heuristics (predictions, no proof)}

\begin{itemize}
\item Heuristic~\ref{heu:joint-density}: density of $n$ with all four $\phi$/$\psi$
simultaneously prime. Standard Hardy--Littlewood with no novel content beyond the local
factors.
\end{itemize}

\subsection*{Open questions / gaps}

\begin{itemize}
\item Question~\ref{q:joint-detection}: can the joint Artin $L$-function over $K$ supply
detection of $\phi_0$-primes beyond what $L(s,\chi_{-4})$ alone gives? We have not found
a positive answer. We suspect ``no'' on orthogonality grounds, but this is not proved.
\item The fundamental-domain bookkeeping in Proposition~\ref{prop:hecke-tree-sum}
is not fully written out; the bijection between primitive Gaussian integers in
$\mathcal{F}$ and first-rows of $\SLNz$-matrices is correct in spirit but the matching of
unit-orbit conventions between the two sides should be checked carefully if one wishes to
extract a precise sum identity (rather than an asymptotic one).
\item Remark~\ref{rmk:positive-side-of-the-coin}: the gap between Heath-Brown narrow
cones (\cite{HB80}) and Landau is quantified as a power $X^{1/2 - 1/4}$ in cone-width.
This is a useful diagnostic but it does not suggest a path to closing the gap.
\item Spin equidistribution on the $b = 1$ slice (Remark~\ref{rmk:spin-on-tree-vs-FI})
is, modulo bookkeeping, equivalent to Conjecture~\ref{conj:bilinear}. We have not
verified this equivalence rigorously; a careful matching of weights is needed.
\item The joint-field discussion in \S\ref{sec:joint} establishes the algebraic structure
but produces no new analytic ammunition. \emph{We consider this section to be primarily
clarificatory; it does not by itself attack Landau.}
\end{itemize}

\subsection*{Net assessment}

The principal contribution of this note is the explicit identification
(Theorem~\ref{thm:Phi0-is-Gaussian-factorization}) of the $\Phi_0$ map with the
Gaussian-integer factorization, made systematic at the level of the Hecke
$L$-functions of $\Q(i)$ via Proposition~\ref{prop:hecke-tree-sum}. This dictionary
allows one to translate between Shakov's combinatorial $\SLNz$-language and the
class-field-theoretic / analytic-number-theoretic language used by Iwaniec, Heath-Brown,
Friedlander--Iwaniec, etc. We also pinpoint the precise sense in which Hecke characters
fail to see Landau (Theorem~\ref{thm:GRH-implies-equidistribution}: arguments of $n+i$
degenerate), and we formulate a single bilinear conjecture
(Conjecture~\ref{conj:bilinear}) on the matrix variables whose proof would yield Landau
via the asymptotic sieve.

\textbf{We do not prove Landau's fourth problem.} We do not improve on Iwaniec's $P_2$
result, on Pascadi's $P^+(n^2+1) \ge n^{1.3}$ greatest-prime-factor bound, or on any
other unconditional result in the literature. What we offer is a re-organisation of the
problem in the language of class field theory of $\Q(i)$, made explicit through the
bijection $\Phi_0$, and a precise localisation of the missing analytic ingredient as a
Type-II bilinear estimate on $\SLN$-words.

\begin{thebibliography}{99}

\bibitem{shakov} A. Shakov,
\emph{Polynomials in $\Z[x]$ whose divisors are enumerated by $\SLNz$},
preprint, 2024, arXiv:2405.03552.

\bibitem{cox} D. A. Cox,
\emph{Primes of the Form $x^2 + ny^2$: Fermat, Class Field Theory, and Complex
Multiplication}, 3rd ed., AMS Chelsea, 2022.

\bibitem{iwaniec1976} H. Iwaniec,
The half-dimensional sieve, \emph{Acta Arithmetica} \textbf{29} (1976), 69--95.

\bibitem{iwaniec1978} H. Iwaniec,
Almost-primes represented by quadratic polynomials,
\emph{Inventiones Math.} \textbf{47} (1978), 171--188.

\bibitem{IK} H. Iwaniec, E. Kowalski,
\emph{Analytic Number Theory},
American Mathematical Society Colloquium Publications \textbf{53}, 2004.

\bibitem{FI98a} J. Friedlander, H. Iwaniec,
Asymptotic sieve for primes,
\emph{Annals of Math.} \textbf{148} (1998), 1041--1065.

\bibitem{FI98b} J. Friedlander, H. Iwaniec,
The polynomial $X^2 + Y^4$ captures its primes,
\emph{Annals of Math.} \textbf{148} (1998), 945--1040.

\bibitem{FI10} J. Friedlander, H. Iwaniec,
\emph{Opera de Cribro},
American Mathematical Society Colloquium Publications \textbf{57}, 2010.

\bibitem{HB80} D. R. Heath-Brown,
The distribution of Gaussian primes,
\emph{Acta Arithmetica}, several papers from 1980 onwards.

\bibitem{hecke1920} E. Hecke,
\emph{Eine neue Art von Zetafunktionen und ihre Beziehungen zur Verteilung der
Primzahlen II},
Mathematische Zeitschrift \textbf{6} (1920), 11--51.

\bibitem{neukirch} J. Neukirch,
\emph{Algebraic Number Theory},
Grundlehren der mathematischen Wissenschaften \textbf{322}, Springer, 1999.

\bibitem{lang-alg} S. Lang,
\emph{Algebra}, 3rd revised ed., Graduate Texts in Mathematics \textbf{211}, Springer,
2002.

\bibitem{R3-friedlander-iwaniec}
Internal note R3, \emph{Friedlander--Iwaniec, parity-breaking, and why $n^2+1$ resists},
this project, \texttt{notes/research/R3-friedlander-iwaniec.md}.

\bibitem{R4-binary-quadratic-forms}
Internal note R4, \emph{Binary quadratic forms, Gauss composition, and Shakov's four
polynomials}, this project, \texttt{notes/research/R4-binary-quadratic-forms.md}.

\end{thebibliography}

\end{document}
